
% =============================================================================
% Project  : THE PHILOSOPHY OF VIABILITY
% Book     : Book I — Ontology of Viability
% Book ID  : VS.B01
% Part     : Part I — Identity and Manifestation
% Part ID  : VS.B01.P01
% Chapter  : Chapter 08 — Comparative Philosophical Examination
% Chapter ID: VS.B01.P01.C08
%
% File     :
%   2_MAINMATTER/PART_I_IDENTITY_AND_MANIFESTATION/
%   CHAPTER_08_COMPARATIVE_PHILOSOPHICAL_EXAMINATION/
%   sections/section_14_suhrawardi.tex
%
% Section  : Section 14 — Suhrawardi
% Section ID: VS.B01.P01.C08.S14
%
% Purpose:
%   Reconstruct and critically examine Suhrawardi's Illuminationist
%   ontology, with particular attention to light, the hierarchy of
%   lights, the relation between essence and existence, self-manifestation,
%   the suspended images, and the ontological question of identity.
%
% Methodological Status:
%   Comparative historical and philosophical examination.
%   No historical doctrine is adopted as a foundational commitment
%   of Book I merely by virtue of its similarity to the present framework.
%
% =============================================================================

\section{Suhrawardi}
\label{sec:ch08_s14}

% -----------------------------------------------------------------------------
% SS01 — Scope and Comparative Orientation
% -----------------------------------------------------------------------------

\subsection{Scope and Comparative Orientation}
\label{sec:ch08-s14-ss01}

% Structural ID: VS.B01.P01.C08.S14.SS01.P001
\phantomsection\label{para:VS.B01.P01.C08.S14.SS01.P001}

The present section examines the philosophy of Shihab al-Din Yahya
al-Suhrawardi, particularly the ontology developed in
\textit{Hikmat al-Ishraq} (\textit{The Philosophy of Illumination}).
The comparison concerns the ontological role assigned to light,
the differentiation of modes of light and darkness, the hierarchy
of immaterial lights, the relation between self-manifestation and
self-awareness, the status of suspended images, and the question
of whether the account of light provides an adequate explanation
of the identity of existent things. The inquiry is restricted to
the problems relevant to Part I of Book I. It does not attempt
to reconstruct the entirety of Suhrawardi's logic, physics,
psychology, political philosophy, or eschatology.

% Structural ID: VS.B01.P01.C08.S14.SS01.P002
\phantomsection\label{para:VS.B01.P01.C08.S14.SS01.P002}

Suhrawardi is relevant because his philosophy does not treat light
merely as a visible property of bodies or as a metaphor for
intelligibility. Light occupies a central ontological position
in his account of reality. The distinctions among pure lights,
accidental lights, dark barriers, and dark accidents are intended
to explain different modes of existence and dependence. His
account therefore raises a question that cannot be settled by
comparing isolated uses of the word ``light'': whether the
ontological differentiation of things through light can also
explain what makes an existent the identity that it is.

% Structural ID: VS.B01.P01.C08.S14.SS01.P003
\phantomsection\label{para:VS.B01.P01.C08.S14.SS01.P003}

The comparison is not based on the assumption that Suhrawardi's
light is equivalent to the Light discussed elsewhere in the
Philosophy of Viability. Nor does it assume that the hierarchy
of lights corresponds to the proposed realms of Book I. The
historical doctrine and the present ontology must first be
reconstructed independently. Any proposed correspondence must
then be assessed according to its conceptual function,
ontological commitments, explanatory scope, and limitations.
The methodological distinction established in
Section~\ref{sec:ch08-method-and-criteria} governs the present
examination.

% Structural ID: VS.B01.P01.C08.S14.SS01.P004
\phantomsection\label{para:VS.B01.P01.C08.S14.SS01.P004}

The central comparative question is therefore not whether
Suhrawardi anticipated the terminology of Book I. It is whether
his account of light, manifestation, and the differentiation
of existents identifies an ontological dimension that the
distinction between Being and Identity must acknowledge, and
whether that account can explain identity without reducing it
to luminosity, intensity, self-awareness, quiddity, or causal
position within a hierarchy.

% -----------------------------------------------------------------------------
% SS02 — Historical and Textual Location
% -----------------------------------------------------------------------------

\subsection{Historical and Textual Location}
\label{sec:ch08-s14-ss02}

% Structural ID: VS.B01.P01.C08.S14.SS02.P001
\phantomsection\label{para:VS.B01.P01.C08.S14.SS02.P001}

Suhrawardi (1154--1191) was trained in the Avicennan Peripatetic
tradition and later developed a philosophical project identified
with \textit{ishraq}. His work cannot be adequately described as
a simple rejection of philosophy in favour of mystical experience.
His major writings address logic, physics, metaphysics, and
cosmology, while his allegorical and Persian writings also
contribute to the reception and interpretation of his thought.
The relation between his Peripatetic writings and
\textit{Hikmat al-Ishraq} must therefore be examined rather than
presumed to constitute a single undifferentiated doctrinal
position \citep{MarcotteSuhrawardiSEP2019,ZiaiKnowledgeIllumination1990}.

% Structural ID: VS.B01.P01.C08.S14.SS02.P002
\phantomsection\label{para:VS.B01.P01.C08.S14.SS02.P002}

The principal primary text for the present comparison is
\textit{Hikmat al-Ishraq}, completed in 1186. Its two principal
discourses must be distinguished: the first develops logical
and methodological criticisms of Peripatetic philosophy, while
the second presents the metaphysics and cosmology of light.
Other writings, including \textit{al-Talwihat}
(\textit{The Intimations}) and \textit{al-Mashari wa-l-Mutarahat}
(\textit{The Paths and Havens}), are relevant where they clarify
the development or qualification of a doctrine. Their statements
must not, however, be transferred indiscriminately into
\textit{Hikmat al-Ishraq} without attention to textual context
and chronology \citep{SuhrawardiPhilosophyIllumination,MarcotteSuhrawardiSEP2019}.

% Structural ID: VS.B01.P01.C08.S14.SS02.P003
\phantomsection\label{para:VS.B01.P01.C08.S14.SS02.P003}

The textual reconstruction also distinguishes Suhrawardi's
explicit formulations from the interpretations of later
commentators. Ibn Kammuna and Shahrzuri, for example, represent
different tendencies in the early reception of his work:
one gives particular emphasis to its discursive and systematic
philosophical dimensions, while the other develops its symbolic
and cosmological implications. Qutb al-Din Shirazi's commentary
subsequently became an influential source for the interpretation
of \textit{Hikmat al-Ishraq}. These commentarial developments
are philosophically important, but they cannot be treated as
direct evidence that every later formulation was already
systematically present in Suhrawardi's own text
\citep{MarcotteSuhrawardiSEP2019,WalbridgeLeavenAncients2000}.

% Structural ID: VS.B01.P01.C08.S14.SS02.P004
\phantomsection\label{para:VS.B01.P01.C08.S14.SS02.P004}

The historical description of Suhrawardi as the founder of a
fully unified Illuminationist school also requires qualification.
The existence of a distinctive philosophical project is not
identical with the existence of a historically continuous,
uniform, and doctrinally homogeneous school. The later
transmission of his writings involved commentaries,
reformulations, selective appropriations, and disagreements.
The present section therefore uses ``Illuminationist'' primarily
to identify the philosophical orientation associated with
Suhrawardi, not to presuppose that every later thinker who
employed his terminology shared one complete doctrine
\citep{MarcotteSuhrawardiSEP2019,VanLitWorldImage2017}.

% Structural ID: VS.B01.P01.C08.S14.SS02.P005
\phantomsection\label{para:VS.B01.P01.C08.S14.SS02.P005}

This textual caution is particularly important for the concept
of the world of images. The later commentarial tradition,
especially the work of Shahrzuri, contributed to the more
systematic formulation and transmission of the expression
``world of image.'' The suspended images discussed by
Suhrawardi must therefore be distinguished from the complete
cosmological system that later commentators constructed around
them. The present examination will refer to the suspended
images and to the additional ontological domain associated
with them, while reserving judgment on the extent to which
the later stabilized terminology can be attributed directly
to Suhrawardi \citep{VanLitWorldImage2017}.

% -----------------------------------------------------------------------------
% SS03 — Textual and Interpretive Disputes
% -----------------------------------------------------------------------------

\subsection{Textual and Interpretive Disputes}
\label{sec:ch08-s14-ss03}

% Structural ID: VS.B01.P01.C08.S14.SS03.P001
\phantomsection\label{para:VS.B01.P01.C08.S14.SS03.P001}

Several interpretive questions affect the reconstruction of
Suhrawardi's ontology. The first concerns the status of light:
whether it should be understood as an ontological reality,
a principle of manifestation, a mode of existence, or a
concept whose function is partly explanatory and partly
phenomenological. The second concerns the relation between
light and substance: whether light itself is a substance,
whether some lights are self-subsistent substances while
others are accidents, or whether the term ``light'' covers
ontologically different kinds of entities. The third concerns
the relation between the suspended images and the later
conceptualization of a distinct world of images.

% Structural ID: VS.B01.P01.C08.S14.SS03.P002
\phantomsection\label{para:VS.B01.P01.C08.S14.SS03.P002}

A further dispute concerns the relation between Suhrawardi's
critique of the distinction between essence and existence
and his positive ontology of light. Interpretations that
describe him as a proponent of the primacy of essence must
be examined alongside his criticism of existence as a
mental consideration and his account of light as an
ontological principle. Neither ``primacy of essence'' nor
``primacy of existence'' should be attributed to him without
specifying the sense in which essence, existence, and light
are being discussed. The issue is not merely terminological:
it affects the account of what is real and the way in which
ontological differentiation is explained
\citep{MarcotteSuhrawardiSEP2019,WalbridgeLeavenAncients2000}.

% Structural ID: VS.B01.P01.C08.S14.SS03.P003
\phantomsection\label{para:VS.B01.P01.C08.S14.SS03.P003}

The historiography of Suhrawardi also contains disagreement
about the extent to which his project represents a radical
departure from Avicennan Peripateticism. One interpretation
emphasizes the recovery of an ancient wisdom that joins Greek
and Iranian traditions. Another emphasizes his critical
reconfiguration of Avicennan philosophy, including its
metaphysics, physics, and theory of definition. These
interpretations are not necessarily mutually exclusive in
every respect, but neither can be accepted as a complete
historical explanation without textual argument. Claims of
direct historical transmission from ancient Greek or Persian
sources must be distinguished from Suhrawardi's own
presentation of an intellectual genealogy
\citep{WalbridgeLeavenAncients2000,ZiaiKnowledgeIllumination1990}.

% Structural ID: VS.B01.P01.C08.S14.SS03.P004
\phantomsection\label{para:VS.B01.P01.C08.S14.SS03.P004}

The present comparison adopts a limited reconstruction.
It accepts the ontological seriousness of Suhrawardi's
metaphysics of light, while refusing to treat every later
Illuminationist interpretation as part of his original
doctrine. It recognizes the additional status assigned to
suspended images, while refusing to assume that their later
cosmological elaboration is identical with the original
argument. It recognizes the importance of knowledge by
presence, while distinguishing that epistemological account
from the ontological independence required by Book I.
These qualifications are not attempts to weaken Suhrawardi's
position. They are conditions for comparing it without
substituting a later or simplified doctrine for the one
actually under examination.

% -----------------------------------------------------------------------------
% SS04 — Illumination and Demonstration
% -----------------------------------------------------------------------------

\subsection{Illumination and Demonstration}
\label{sec:ch08-s14-ss04}

% Structural ID: VS.B01.P01.C08.S14.SS04.P001
\phantomsection\label{para:VS.B01.P01.C08.S14.SS04.P001}

The term \textit{ishraq} is associated with illumination,
radiance, and the disclosure of light. In Suhrawardi's
philosophical project, illumination is not merely a literary
image attached to an otherwise Peripatetic system. It names
a mode of access to reality that is connected with the
ontological centrality of light. The philosophical method
therefore includes both discursive demonstration and a
form of direct apprehension that Suhrawardi regards as
indispensable for the realization of the truths discussed
in his metaphysics \citep{ZiaiKnowledgeIllumination1990,MarcotteSuhrawardiSEP2019}.

% Structural ID: VS.B01.P01.C08.S14.SS04.P002
\phantomsection\label{para:VS.B01.P01.C08.S14.SS04.P002}

The relation between demonstration and illumination must
not be reduced to a simple opposition between reason and
experience. Suhrawardi criticizes the sufficiency of
Peripatetic definition and discursive reasoning for certain
metaphysical questions, but he also employs arguments,
distinctions, and demonstrations throughout his writings.
His position is therefore not adequately represented by
the claim that reasoning has no philosophical value.
Rather, he disputes the claim that discursive definition
and inference alone exhaust the possible access to
ontological reality \citep{SuhrawardiPhilosophyIllumination,ZiaiKnowledgeIllumination1990}.

% Structural ID: VS.B01.P01.C08.S14.SS04.P003
\phantomsection\label{para:VS.B01.P01.C08.S14.SS04.P003}

The doctrine of the master of the two orders concerns the
conjunction of philosophical competence and illuminative
realization. The ideal philosopher is not merely a
technician of demonstrations, nor is the ideal sage merely
a person who reports private experiences. The ideal
requires the conjunction of discursive and intuitive
capacities. This doctrine is relevant to the historical
and methodological character of Suhrawardi's project,
but it does not by itself establish an ontological category
or a principle of identity.

% Structural ID: VS.B01.P01.C08.S14.SS04.P004
\phantomsection\label{para:VS.B01.P01.C08.S14.SS04.P004}

Book I does not adopt illumination, mystical intuition,
or any other cognitive act as a condition of the existence
of Being or Identity. The comparison can acknowledge that
Suhrawardi regards direct apprehension as indispensable
to his philosophical method without accepting that
ontological reality depends upon being apprehended.
The question of how a knower accesses reality belongs
to a different inquiry from the question of what exists
independently of the knower. The distinction is essential
to prevent the epistemological dimension of Suhrawardi's
philosophy from being imported into the ontological
foundations of this volume.

% -----------------------------------------------------------------------------
% SS05 — Light as Ontological Principle
% -----------------------------------------------------------------------------

\subsection{Light as an Ontological Principle}
\label{sec:ch08-s14-ss05}

% Structural ID: VS.B01.P01.C08.S14.SS05.P001
\phantomsection\label{para:VS.B01.P01.C08.S14.SS05.P001}

Light occupies a fundamental position in Suhrawardi's
account of reality. It is not simply one visible quality
among others, nor is its significance exhausted by its
ordinary physical occurrence. Light is self-evident in
the sense that it does not require a more familiar concept
to make its reality intelligible. The argument begins
from the claim that light is manifest and that its
manifestness is not adequately explained by treating it
as an ordinary object whose nature is accessible only
through an external definition \citep[Part II, §§107--115]{SuhrawardiPhilosophyIllumination}.

% Structural ID: VS.B01.P01.C08.S14.SS05.P002
\phantomsection\label{para:VS.B01.P01.C08.S14.SS05.P002}

The self-evidence of light must be distinguished from
the claim that every use of the word ``light'' designates
one simple kind of substance. Suhrawardi's ontology
contains self-subsistent lights, accidental lights,
dark substances, and dark accidents. The term therefore
operates across distinctions of dependence and mode
of existence. A statement about the self-evidence of
light cannot be converted without further argument
into the universal proposition that everything real
is a substance of one homogeneous luminous kind
\citep{WalbridgeLeavenAncients2000,MarcotteSuhrawardiSEP2019}.

% Structural ID: VS.B01.P01.C08.S14.SS05.P003
\phantomsection\label{para:VS.B01.P01.C08.S14.SS05.P003}

Light is also connected with manifestation. Pure lights
are manifest to themselves and, under appropriate
relations, manifest to other lights. Accidental light
makes an otherwise dark or non-self-manifesting entity
manifest in a different manner. The ontological account
therefore connects luminosity with degrees and modes
of manifestation. Yet manifestation, in this context,
must not be equated with the concept of Manifestation
developed in Book I. The shared word identifies a
possible comparative problem, not an established
identity of conceptual function.

% Structural ID: VS.B01.P01.C08.S14.SS05.P004
\phantomsection\label{para:VS.B01.P01.C08.S14.SS05.P004}

The claim that light is fundamental does not by itself
establish that light is identical with Being, Identity,
Determination, Unity, or any one of the proposed realms
of Book I. Nor does the self-manifesting character of
light establish that identity is constituted by
self-manifestation. These conclusions would require
additional arguments concerning the relation between
ontological reality, manifestation, and the identity
of the existent. Suhrawardi's account makes that relation
a serious problem; it does not settle it merely by
assigning a central role to light.

% Structural ID: VS.B01.P01.C08.S14.SS05.P005
\phantomsection\label{para:VS.B01.P01.C08.S14.SS05.P005}

The ontological status of Light in Book I consequently
remains open to its designated investigation. The
present section neither adopts Suhrawardi's metaphysics
as the definition of Light nor rejects it by stipulating
that Light must be a process, a relation, a substance,
or a structural medium. The comparison establishes
that any later account of Light must be capable of
addressing the distinctions between self-subsistence,
dependence, manifestation, and the different kinds
of luminous and dark entities.

% -----------------------------------------------------------------------------
% SS06 — Light and Darkness
% -----------------------------------------------------------------------------

\subsection{Light and Darkness}
\label{sec:ch08-s14-ss06}

% Structural ID: VS.B01.P01.C08.S14.SS06.P001
\phantomsection\label{para:VS.B01.P01.C08.S14.SS06.P001}

Suhrawardi's ontology does not treat the distinction
between light and darkness as a simple contrast between
the presence and absence of a visible quality. Darkness
has an articulated role in the classification of
existent things. Dark barriers and dark accidents
are not merely verbal negations of luminous entities.
They belong to the account of the material and
dependent modes of reality. This distinction prevents
the metaphysics of light from being reduced to the
claim that whatever is not luminous is simply
nothing \citep[Part II, §§107--113]{SuhrawardiPhilosophyIllumination}.

% Structural ID: VS.B01.P01.C08.S14.SS06.P002
\phantomsection\label{para:VS.B01.P01.C08.S14.SS06.P002}

The relation between light and darkness is nevertheless
asymmetrical. Pure light has a mode of manifestation
that dark barriers do not possess in the same manner.
Darkness is associated with the absence of self-manifestation
and with dependence upon another for manifestation.
This asymmetry is central to Suhrawardi's explanatory
order, but it does not warrant the conclusion that
dark entities are unreal or that their ontological
status is exhausted by a privation. Their place in
the classification must be reconstructed through
the distinctions among substance, accident, and
dependence.

% Structural ID: VS.B01.P01.C08.S14.SS06.P003
\phantomsection\label{para:VS.B01.P01.C08.S14.SS06.P003}

The earlier versions of the Philosophy of Viability
sometimes described Suhrawardi's darkness as a pure
privation and contrasted it with a positive concept
of sediment. That comparison is too compressed.
It risks treating every dark entity in Suhrawardi's
ontology as a mere absence, while simultaneously
using ``sediment'' as if it were an established
equivalent. The present examination does not adopt
that equation. The relation between darkness,
privation, material constitution, and any account
of accumulated effects must be examined separately.

% Structural ID: VS.B01.P01.C08.S14.SS06.P004
\phantomsection\label{para:VS.B01.P01.C08.S14.SS06.P004}

The distinction between light and darkness also
does not map directly onto the distinction between
identity and non-identity. A dark entity is not
thereby an entity without identity. Nor does the
degree of luminosity, by itself, establish whether
an existent is numerically one, structurally unified,
individuated, or identical with itself. These are
different ontological questions. Any relation among
them must be established rather than inferred from
the contrast between light and darkness.

% -----------------------------------------------------------------------------
% SS07 — Four Classes of Existent
% -----------------------------------------------------------------------------

\subsection{The Four Classes of Existent}
\label{sec:ch08-s14-ss07}

% Structural ID: VS.B01.P01.C08.S14.SS07.P001
\phantomsection\label{para:VS.B01.P01.C08.S14.SS07.P001}

Suhrawardi's classification of existent things distinguishes
four principal classes: pure immaterial lights, dark
barriers, accidental lights, and dark accidents. These
classes are not four equivalent degrees of one material
continuum. They distinguish different modes of
self-subsistence, inherence, manifestation, and
dependence. Their function is to articulate the
ontological structure of the world rather than merely
to provide a vocabulary for describing visible
differences \citep[Part II, §§107--113, 121]{SuhrawardiPhilosophyIllumination}.

% Structural ID: VS.B01.P01.C08.S14.SS07.P002
\phantomsection\label{para:VS.B01.P01.C08.S14.SS07.P002}

Pure immaterial lights are self-subsistent and
incorporeal. Their luminosity is not an accidental
quality inhering in a material body. They are
manifest to themselves and possess a mode of
awareness that is not dependent upon the reception
of a physical form. They occupy the highest
ontological position among the classes of beings
other than the Light of Lights. Their self-subsistence
must not, however, be translated into the claim
that they are identities in the sense developed
in Book I.

% Structural ID: VS.B01.P01.C08.S14.SS07.P003
\phantomsection\label{para:VS.B01.P01.C08.S14.SS07.P003}

Dark barriers are corporeal entities that do not
manifest themselves in the manner of pure lights.
They are associated with extension, bodily
constitution, and the reception of properties.
Their ontological status is not reducible to the
absence of a visible light source. They form the
material dimension of Suhrawardi's cosmology,
although their existence and manifestation are
explained through relations with incorporeal and
accidental lights.

% Structural ID: VS.B01.P01.C08.S14.SS07.P004
\phantomsection\label{para:VS.B01.P01.C08.S14.SS07.P004}

Accidental lights are luminous entities or states
that inhere in another. They differ from pure
immaterial lights because they do not possess
self-subsistence in the same manner. Their
manifestation is dependent upon a subject in
which they occur. Dark accidents, by contrast,
are non-luminous qualities or states that also
inhere in another. The distinction between
accidental light and dark accident is therefore
not simply a distinction between two substances;
it is a distinction within the domain of dependent
states and qualities.

% Structural ID: VS.B01.P01.C08.S14.SS07.P005
\phantomsection\label{para:VS.B01.P01.C08.S14.SS07.P005}

The fourfold classification must not be mapped
onto the four fundamental realms proposed in
Book I. The two classifications answer different
questions. Suhrawardi's classes distinguish
modes of subsistence, inherence, and manifestation
within his ontology of light and darkness. The
four realms of Book I are proposed as distinct
ontological domains of determination and
manifestation. Similarity in the number of
classes or the use of a term such as ``realm''
would not establish a common structure or
a one-to-one correspondence.

% Structural ID: VS.B01.P01.C08.S14.SS07.P006
\phantomsection\label{para:VS.B01.P01.C08.SS07.P006}

The four classes also do not provide a complete
theory of Identity. They describe ontological
differences among existent things, but they do
not by themselves establish what makes one
pure light distinct from another, whether a
dark barrier can possess more than one identity,
or whether the identity of an entity is exhausted
by its mode of subsistence. Those questions
remain open within the comparative inquiry.

% -----------------------------------------------------------------------------
% SS08 — Light of Lights and Causal Order
% -----------------------------------------------------------------------------

\subsection{The Light of Lights and the Causal Order}
\label{sec:ch08-s14-ss08}

% Structural ID: VS.B01.P01.C08.S14.SS08.P001
\phantomsection\label{para:VS.B01.P01.C08.S14.SS08.P001}

The Light of Lights (\textit{nur al-anwar}) occupies the
highest position in Suhrawardi's hierarchy. It is
the first principle from which the order of lights
proceeds. Its primacy is not merely a matter of
being the brightest entity in a series of visible
objects. It is an ontological and causal primacy
within the structure of reality. Suhrawardi's
arguments for this principle draw upon the
principle of nobler contingency and arguments
concerning the impossibility of an ordered,
actual infinite regress in the relevant causal
series \citep[Part II, §§129--138, 164--173]{SuhrawardiPhilosophyIllumination}.

% Structural ID: VS.B01.P01.C08.S14.SS08.P002
\phantomsection\label{para:VS.B01.P01.C08.S14.SS08.P002}

The Light of Lights is identified in Suhrawardi's
system with the divine first principle. The
dependent lights derive their ontological
position from their relation to this principle,
while remaining distinct within the hierarchy.
The relation between the first principle and
the subsequent lights is therefore not a relation
of numerical identity. Nor does the fact that
all dependent lights derive from one source
establish that they constitute one identity or
that their plurality is merely apparent.

% Structural ID: VS.B01.P01.C08.S14.SS08.P003
\phantomsection\label{para:VS.B01.P01.C08.S14.SS08.P003}

The Light of Lights must not be identified with
Being as such, with Identity, or with First
Determination in the vocabulary of Book I.
The comparison may identify a structural
question concerning first principles and
ontological dependence, but the concepts
must not be substituted for one another.
The Light of Lights has a specific function
within Suhrawardi's theological and cosmological
argument. First Determination has a different
proposed function within Book I. Their
equivalence cannot be established by their
respective positions near the beginning of
an explanatory order.

% Structural ID: VS.B01.P01.C08.S14.SS08.P004
\phantomsection\label{para:VS.B01.P01.C08.S14.SS08.P004}

The hierarchy also raises a question concerning
the relation between causal dependence and
ontological differentiation. A dependent light
may receive its existence or its mode of
manifestation from a higher principle without
being reducible to that principle. The causal
origin of an existent does not, by itself,
explain its identity. A theory of causal
dependence and a theory of identity may
intersect, but neither can be substituted
for the other without an additional argument.

% -----------------------------------------------------------------------------
% SS09 — Vertical and Horizontal Orders
% -----------------------------------------------------------------------------

\subsection{The Vertical and Horizontal Orders of Lights}
\label{sec:ch08-s14-ss09}

% Structural ID: VS.B01.P01.C08.S14.SS09.P001
\phantomsection\label{para:VS.B01.P01.C08.S14.SS09.P001}

Suhrawardi's hierarchy is not adequately represented
as one simple linear sequence of lights. He
distinguishes a vertical order proceeding from
the Light of Lights and a horizontal order
involving lights that possess equal or comparable
ontological rank. The vertical order includes
the dominating lights, while the horizontal
order includes the mother-lights and the
Lords of Species associated with the
intelligible forms of natural kinds. These
orders interact in the generation and
differentiation of the lower world
\citep[Part II, §§141--156, 183--184]{SuhrawardiPhilosophyIllumination}.

% Structural ID: VS.B01.P01.C08.S14.SS09.P002
\phantomsection\label{para:VS.B01.P01.C08.S14.SS09.P002}

The vertical order concerns ontological dependence
and the relation of higher to lower lights.
The horizontal order is not simply another
stage in the same descending chain. It
introduces a different form of plurality
among immaterial lights, including lights
that may be equal in rank while differing
in their aspects and relations. The two
orders therefore cannot be reduced to
one numerical sequence without losing
the distinction that Suhrawardi's system
attempts to preserve.

% Structural ID: VS.B01.P01.C08.S14.SS09.P003
\phantomsection\label{para:VS.B01.P01.C08.S14.SS09.P003}

The distinction between the vertical and
horizontal orders is relevant to the
problem of plurality. It shows that
ontological differentiation is not
explained solely by the position of an
entity in a single hierarchy. Equal
rank does not entail numerical identity,
and difference of rank does not by
itself establish a difference of identity.
The system therefore raises a problem
that remains open in Book I: whether
plurality of identities can be adequately
explained through a plurality of
ontological positions, aspects, or
relations.

% Structural ID: VS.B01.P01.C08.S14.SS09.P004
\phantomsection\label{para:VS.B01.P01.C08.S14.SS09.P004}

The following diagram is a schematic reconstruction
of the two orders and the four principal classes
of existent things. It is not a temporal sequence,
a complete cosmological map, or a proposed
correspondence with the four fundamental
realms and three intermediate hubs of Book I.


% =============================================================================
% Project : THE PHILOSOPHY OF VIABILITY
% Book    : Book I — Ontology of Viability
% Book ID : VS.B01
% Chapter : Chapter 08 — Comparative Philosophical Examination
% Section : Section 14 — Suhrawardi
% Asset   : Diagram 01 — Orders of Light and Classes of Existent Things
% Asset ID: VS.B01.P01.C08.S14.SS09.P004.DG01
% File    : 4_ASSETS/diagrams/VS.B01.P01.C08.S14.SS09.P004.DG01.suhrawardi_orders_of_light.tex
%
% Qualification:
%   The two panels represent distinct classificatory dimensions.
%   The four classes on the right are not arranged as a causal sequence.
% =============================================================================

\begin{figure}[htbp]
\centering

\begin{tikzpicture}[
    >=Stealth,
    font=\small,
    lightnode/.style={
        draw=ViabilityDarkBlue,
        fill=ViabilityPaleBlue,
        rounded corners=2pt,
        align=center,
        text width=5.0cm,
        minimum height=10mm,
        inner sep=5pt
    },
    ordernode/.style={
        draw=ViabilityGold,
        fill=ViabilityPaleGold,
        rounded corners=2pt,
        align=center,
        text width=4.8cm,
        minimum height=12mm,
        inner sep=5pt
    },
    classnode/.style={
        draw=ViabilityGreen,
        fill=ViabilityPaleGreen,
        rounded corners=2pt,
        align=center,
        text width=5.0cm,
        minimum height=11mm,
        inner sep=5pt
    },
    paneltitle/.style={
        font=\bfseries,
        align=center,
        text width=5.2cm
    }
]

% Panel A: Orders of lights

\node[paneltitle] at (0,0)
    {Two Orders of Immaterial Lights};

\node[lightnode] (lol) at (0,-1.0)
    {\textbf{Light of Lights}\\
     First principle of the Illuminationist order};

\node[ordernode] (vertical) at (-1.45,-2.7)
    {\textbf{Vertical Order}\\
     Dominating lights and ontological dependence};

\node[ordernode] (horizontal) at (1.45,-2.7)
    {\textbf{Horizontal Order}\\
     Equal-rank lights, mother-lights, and Lords of Species};

\node[lightnode] (managing) at (0,-4.45)
    {\textbf{Managing Lights}\\
     Soul-related lights associated with corporeal orders};

\draw[->,thick,draw=ViabilityDarkBlue]
    (lol.south) -- (vertical.north);

\draw[<->,dashed,thick,draw=ViabilityGold]
    (vertical.east) -- (horizontal.west);

\draw[densely dotted,thick,draw=ViabilityDarkBlue]
    (vertical.south) -- (managing.north);

\node[align=center,font=\scriptsize,text width=5.2cm]
    at (0,-5.35)
    {The dotted connection indicates association, not a further
     stage in a single linear sequence};

% Panel B: Four classes

\node[paneltitle] at (8.2,0)
    {Four Principal Classes of Existent Things};

\node[classnode] at (8.2,-1.0)
    {\textbf{Pure Immaterial Lights}\\
     Self-subsistent and incorporeal};

\node[classnode] at (8.2,-2.2)
    {\textbf{Dark Barriers}\\
     Corporeal and extended entities};

\node[classnode] at (8.2,-3.4)
    {\textbf{Accidental Lights}\\
     Luminous states inhering in another};

\node[classnode] at (8.2,-4.6)
    {\textbf{Dark Accidents}\\
     Non-luminous states inhering in another};

\node[
    draw=gray,
    fill=white,
    rounded corners=2pt,
    align=center,
    text width=11.0cm,
    inner sep=5pt
] at (4.1,-6.0)
    {\textit{The two panels show different classificatory dimensions.
    No one-to-one correspondence between them is implied.}};

\end{tikzpicture}

\caption{Schematic distinction between Suhrawardi's vertical and
horizontal orders of immaterial lights and his four principal
classes of existent things. The diagram is non-temporal and
does not imply that the four classes form a causal sequence.}
\label{fig:ch08-s14-suhrawardi-orders-of-light}

\end{figure}

% =============================================================================
% END OF DIAGRAM 01 — SUHRAWARDI
% =============================================================================


% -----------------------------------------------------------------------------
% SS10 — Dominance, Love, and Relation
% -----------------------------------------------------------------------------

\subsection{Dominance, Love, and Relation}
\label{sec:ch08-s14-ss10}

% Structural ID: VS.B01.P01.C08.S14.SS10.P001
\phantomsection\label{para:VS.B01.P01.C08.S14.SS10.P001}

The relations among Suhrawardi's lights include
dominance (\textit{qahr}) and love or yearning
(\textit{mahabbah} and related terms). Higher
lights exercise a form of domination over
lower lights, while lower lights desire or
yearn toward higher ones. These relations
are not merely psychological descriptions
of human emotion. They articulate an
ontological order in which causation,
dependence, perfection, and orientation
toward a higher principle are interconnected
\citep[Part II, §§147--156]{SuhrawardiPhilosophyIllumination}.

% Structural ID: VS.B01.P01.C08.S14.SS10.P002
\phantomsection\label{para:VS.B01.P01.C08.S14.SS10.P002}

The relation of dominance does not mean that
a lower light is simply identical with the
higher light that causes it. Nor does
yearning establish that the lower light
is merely a deficient part of the higher.
The relations are asymmetric, but the
entities related through them remain
distinguishable within the hierarchy.
The system thus provides a relational
account of ontological order without
thereby proving that relation is the
sole ground of identity.

% Structural ID: VS.B01.P01.C08.S14.SS10.P003
\phantomsection\label{para:VS.B01.P01.C08.S14.SS10.P003}

The comparison with Book I must avoid
equating Suhrawardi's love or yearning
with the concept of viability, with
the movement of manifestation toward
a source, or with any proposed rule
of a Book I hub. The relation may
provide a comparative question about
orientation, dependence, and the
differentiation of levels, but its
theological and cosmological role
is specific to Suhrawardi's system.

% Structural ID: VS.B01.P01.C08.S14.SS10.P004
\phantomsection\label{para:VS.B01.P01.C08.S14.SS10.P004}

The two orders of light and their relations
also complicate any simple account of
causation as a one-directional chain.
Multiple lights, including lights of
equal rank, participate in the
generation and differentiation of
lower realities. The resulting order
is structurally richer than a single
linear emanation sequence. Whether
this model provides an adequate
explanation of plurality, individuation,
and identity remains a distinct
question and will be examined below.

% -----------------------------------------------------------------------------
% SS11 — Essence, Existence, and Quiddity
% -----------------------------------------------------------------------------

\subsection{Essence, Existence, and Quiddity}
\label{sec:ch08-s14-ss11}

% Structural ID: VS.B01.P01.C08.S14.SS11.P001
\phantomsection\label{para:VS.B01.P01.C08.S14.SS11.P001}

Suhrawardi's critique of the Peripatetic distinction
between essence and existence is one of the
most consequential parts of his metaphysics.
He challenges the treatment of existence as
a real predicate added to quiddity in the
external world. In the relevant discussion,
existence and certain related notions are
treated as mental considerations rather than
as independently existing realities
\citep[Part II, §§56--68]{SuhrawardiPhilosophyIllumination}.

% Structural ID: VS.B01.P01.C08.S14.SS11.P002
\phantomsection\label{para:VS.B01.P01.C08.S14.SS11.P002}

This critique must not be simplified into the
claim that Suhrawardi merely replaces
existence with essence. His account of
light introduces a different ontological
principle, while the status of quiddity
and the conceptual treatment of existence
remain subjects of interpretive dispute.
The relation between the concept of
existence and the reality of what exists
must be distinguished from the relation
between essence and the concrete
differentiation of existent things.

% Structural ID: VS.B01.P01.C08.S14.SS11.P003
\phantomsection\label{para:VS.B01.P01.C08.S14.SS11.P003}

The present comparison must preserve the
difference between Suhrawardi's
essence--existence problem and the
Being--Identity distinction of Book I.
The two distinctions are not equivalent.
Book I does not treat Identity as
quiddity, and it does not treat Being
as an external existence-predicate
added to an independently constituted
essence. The analytical distinction
between Being and Identity is maintained
together with their ontological
inseparability and the absence of
ontological priority between them.

% Structural ID: VS.B01.P01.C08.S14.SS11.P004
\phantomsection\label{para:VS.B01.P01.C08.S14.SS11.P004}

Suhrawardi nevertheless raises a serious
challenge. If the conceptual distinction
between essence and existence does not
correspond to two independently
existing constituents, then any ontology
that introduces a distinction between
Being and Identity must explain why
that distinction is not merely a
distinction of thought. Book I cannot
answer this challenge by repeating
that Being and Identity are analytically
distinct. It must also explain the
ontological significance of the
distinction while preserving their
inseparability.

% Structural ID: VS.B01.P01.C08.S14.SS11.P005
\phantomsection\label{para:VS.B01.P01.C08.S14.SS11.P005}

The strongest comparative result is therefore
not that Suhrawardi refutes the
Being--Identity distinction, nor that
his critique confirms it. Rather,
his treatment of existence, essence,
and light requires Book I to clarify
the difference between a conceptual
distinction, an ontological distinction,
and an ontologically inseparable
analytical distinction. The precise
grounds for this clarification remain
a task for the positive account of
Identity developed in Part I.

% -----------------------------------------------------------------------------
% SS12 — Self-Awareness and Knowledge by Presence
% -----------------------------------------------------------------------------

\subsection{Self-Awareness and Knowledge by Presence}
\label{sec:ch08-s14-ss12}

% Structural ID: VS.B01.P01.C08.S14.SS12.P001
\phantomsection\label{para:VS.B01.P01.C08.S14.SS12.P001}

Suhrawardi's account of pure light is closely
connected with self-awareness and knowledge
by presence (\textit{ilm huduri}). Pure
immaterial lights are manifest to themselves,
and their self-awareness is not explained
as the reception of an external representation
or the acquisition of an image of themselves.
The soul's awareness of itself provides
an important model for this account of
unmediated presence \citep[Part II, §§114--117]{SuhrawardiPhilosophyIllumination,ZiaiKnowledgeIllumination1990}.

% Structural ID: VS.B01.P01.C08.S14.SS12.P002
\phantomsection\label{para:VS.B01.P01.C08.S14.SS12.P002}

The doctrine of presence is important because
it challenges the assumption that every
relation of awareness must be mediated
through a representation that stands
between the knower and the known.
Suhrawardi uses the self-awareness of
the soul and the manifestness of light
to develop an account of immediate
apprehension. The account has
epistemological consequences, but it
also presupposes an ontological
conception of the kind of entity
capable of such self-manifestation.

% Structural ID: VS.B01.P01.C08.S14.SS12.P003
\phantomsection\label{para:VS.B01.P01.C08.S14.SS12.P003}

Book I must distinguish the ontological
claim that an existent may be manifest
to itself from the epistemological
claim that a subject knows itself
through an unmediated act of awareness.
The former may be relevant to an
account of the mode of existence of
certain entities. The latter concerns
the conditions and character of
knowledge. Neither can be substituted
for the other without argument.

% Structural ID: VS.B01.P01.C08.S14.SS12.P004
\phantomsection\label{para:VS.B01.P01.C08.S14.SS12.P004}

The ontological independence of Identity
from the knower remains a foundational
commitment of Book I. Consequently,
self-awareness cannot be treated as
a universal condition for the existence
of identity. An entity does not acquire
its identity merely because it becomes
aware of itself, and the absence of
self-awareness does not by itself
establish the absence of identity.
The possibility that certain kinds
of beings possess self-awareness is
compatible with the existence of
entities whose identity does not
depend upon any cognitive act.

% Structural ID: VS.B01.P01.C08.S14.SS12.P005
\phantomsection\label{para:VS.B01.P01.C08.S14.SS12.P005}

The comparison therefore does not
import knowledge by presence into
the ontological foundations of
Book I. It records a genuine
difference between Suhrawardi's
account of self-manifestation and
the present inquiry into Identity.
The further philosophical treatment
of knowledge by presence, representation,
and the conditions of cognition
belongs to the epistemological
investigation and must not be
presupposed here.

% -----------------------------------------------------------------------------
% SS13 — Suspended Images
% -----------------------------------------------------------------------------

\subsection{Suspended Images and the Additional Ontological Domain}
\label{sec:ch08-s14-ss13}

% Structural ID: VS.B01.P01.C08.S14.SS13.P001
\phantomsection\label{para:VS.B01.P01.C08.S14.SS13.P001}

Suhrawardi's discussion of suspended images
(\textit{suwar muallaqa}) introduces an
additional ontological problem. These
images are not simply material forms
in a physical subject, nor are they
identical with the pure intelligible
lights. They are described as forms
that possess a mode of existence
independent of ordinary material
inherence. Their status is connected
with the imagination, visionary
experience, and the possibility
of posthumous perception
\citep[Part II, §§244--259]{SuhrawardiPhilosophyIllumination,MarcotteSuhrawardiSEP2019}.

% Structural ID: VS.B01.P01.C08.S14.SS13.P002
\phantomsection\label{para:VS.B01.P01.C08.S14.SS13.P002}

The suspended images must not be reduced
to private mental representations.
Suhrawardi's argument assigns them
a status that is not exhausted by
their occurrence as psychological
events within an individual subject.
Their independence from ordinary
material inherence is central to
the explanation of their role in
vision, imagination, and eschatology.
The ontological question is therefore
whether forms can possess a reality
that is neither the reality of
material bodies nor that of
immaterial intellects.

% Structural ID: VS.B01.P01.C08.S14.SS13.P003
\phantomsection\label{para:VS.B01.P01.C08.S14.SS13.P003}

The later notion of a fully articulated
world of images must nevertheless
be distinguished from the original
argument concerning suspended forms.
The commentarial tradition developed
the terminology, location, and
cosmological implications of this
additional domain. The distinction
between Suhrawardi's own formulations
and the later systematic interpretation
is necessary to avoid attributing
to him a completed cosmology that
may depend upon subsequent
commentarial elaboration
\citep{VanLitWorldImage2017}.

% Structural ID: VS.B01.P01.C08.S14.SS13.P004
\phantomsection\label{para:VS.B01.P01.C08.S14.SS13.P004}

The suspended images raise a problem
for any ontology that assumes that
every manifestation must occur
either in a material body or in
an immaterial intellect. They
suggest that the range of possible
modes of manifestation may be
more complex. This challenge is
relevant to Book I because the
ontology of manifestation must
not be restricted in advance to
the categories already proposed.
It does not, however, establish
that the suspended images belong
to the Hub of Signification or
to Relational Manifestation.

% Structural ID: VS.B01.P01.C08.S14.SS13.P005
\phantomsection\label{para:VS.B01.P01.C08.S14.SS13.P005}

The ontological status of the suspended
images remains a comparative
question rather than an adopted
commitment. Their existence,
independence, and mode of
manifestation must be examined
on the basis of Suhrawardi's
arguments and the subsequent
commentarial tradition. Book I
must not resolve the question
by translating the images into
its own terminology before the
relevant ontological distinctions
have been established.

% -----------------------------------------------------------------------------
% SS14 — Bodies and Barzakh
% -----------------------------------------------------------------------------

\subsection{Bodies, Barriers, and the Meaning of Barzakh}
\label{sec:ch08-s14-ss14}

% Structural ID: VS.B01.P01.C08.S14.SS14.P001
\phantomsection\label{para:VS.B01.P01.C08.S14.SS14.P001}

The term \textit{barzakh} requires particular
care in the reconstruction of Suhrawardi.
In the relevant cosmological discussions,
bodies are described as barriers or
intermediate entities that are not
self-manifesting in the manner of
pure lights. The term may also occur
in descriptions of celestial bodies
and the cosmological order. It must
not be assumed that every occurrence
of \textit{barzakh} has the same
technical function or that it already
expresses the later metaphysical
account associated with Ibn Arabi.

% Structural ID: VS.B01.P01.C08.S14.SS14.P002
\phantomsection\label{para:VS.B01.P01.C08.S14.SS14.P002}

Suhrawardi's use of the term must also
be distinguished from the later
conceptual elaboration of the
intermediate world of images.
A corporeal barrier and a suspended
image are not identical merely
because both may be discussed
through language of intermediacy.
The former belongs to the account
of bodily and material entities;
the latter concerns forms that
do not inhere in ordinary matter
in the same manner. Their relation
must be established through
the specific argument in which
each appears.

% Structural ID: VS.B01.P01.C08.S14.SS14.P003
\phantomsection\label{para:VS.B01.P01.C08.S14.SS14.P003}

The term \textit{barzakh} must not be
equated with the three intermediate
hubs proposed in Book I. The
similarity of terminology does
not establish identity of function.
The Hub of Decree, the Hub of
Signification, and the Hub of
Realization are not introduced
as corporeal barriers, suspended
images, or cosmological spheres.
Any comparison would require
an independent account of the
respective ontological roles
and the conditions under which
the comparison is valid.

% Structural ID: VS.B01.P01.C08.S14.SS14.P004
\phantomsection\label{para:VS.B01.P01.C08.S14.SS14.P004}

The same restriction applies to the
proposed realms of Book I. Neither
the corporeal barriers nor the
world of suspended images can
be assigned to a Book I realm
solely on the basis of their
intermediate position in
Suhrawardi's cosmology. Their
ontological functions must first
be established within his own
system. The comparative result
may remain partial where the
available concepts do not
support a defensible translation.

% -----------------------------------------------------------------------------
% SS15 — Individuation and Identity
% -----------------------------------------------------------------------------

\subsection{Individuation, Self-Manifestation, and Identity}
\label{sec:ch08-s14-ss15}

% Structural ID: VS.B01.P01.C08.S14.SS15.P001
\phantomsection\label{para:VS.B01.P01.C08.S14.SS15.P001}

Suhrawardi's account of individual lights
raises a question concerning the
principle by which one light is
distinguished from another. Differences
of rank, intensity, relation, and
luminous or dark accidents contribute
to the differentiation of entities
within the hierarchy. Yet these
features must not be assumed to
be identical with the identity
of the entities they distinguish.
A theory of individuation and
a theory of Identity answer
related but non-identical questions.

% Structural ID: VS.B01.P01.C08.S14.SS15.P002
\phantomsection\label{para:VS.B01.P01.C08.S14.SS15.P002}

Self-manifestation also cannot be
equated with Identity. A pure
light's awareness of itself may
characterize its mode of existence,
but it does not follow that
identity is an act of awareness
or that every entity must
manifest itself to itself in
order to be what it is. The
ontological independence of
Identity from cognition prevents
the conversion of self-awareness
into a universal criterion
of identity.

% Structural ID: VS.B01.P01.C08.S14.SS15.P003
\phantomsection\label{para:VS.B01.P01.C08.S14.SS15.P003}

Nor can luminosity be treated as
a substitute for individuation.
Two lights may differ in intensity
without the difference of intensity
being sufficient to explain
their numerical distinction.
Conversely, two entities may
share a general mode of luminosity
without thereby becoming one
entity. The relation among
intensity, distinction, and
identity must therefore be
examined rather than assumed.

% Structural ID: VS.B01.P01.C08.S14.SS15.P004
\phantomsection\label{para:VS.B01.P01.C08.S14.SS15.P004}

Suhrawardi's account thus places
pressure on Book I to explain
how Identity is positively
distinguished from the conditions
and features through which
an existent is manifested.
Book I cannot answer this
problem by equating Identity
with determination, structure,
luminosity, self-awareness,
or a signature of differences.
It must explain the ontological
role of Identity without
reducing it to any one of
these differentiating features.

% -----------------------------------------------------------------------------
% SS16 — Points of Convergence and Divergence
% -----------------------------------------------------------------------------

\subsection{Points of Convergence and Divergence}
\label{sec:ch08-s14-ss16}

% Structural ID: VS.B01.P01.C08.S14.SS16.P001
\phantomsection\label{para:VS.B01.P01.C08.S14.SS16.P001}

The first point of convergence concerns
the refusal to reduce ontological
reality to its immediate material
appearance. Suhrawardi's ontology
distinguishes pure lights, accidental
lights, dark barriers, and suspended
images. Book I likewise does not
assume that the material manifestation
of an existent exhausts its ontological
status. This convergence concerns
the recognition of ontological
complexity; it does not establish
that the two systems identify
the same domains or entities.

% Structural ID: VS.B01.P01.C08.S14.SS16.P002
\phantomsection\label{para:VS.B01.P01.C08.S14.SS16.P002}

The second point of convergence concerns
the irreducibility of manifestation
to a merely verbal description.
In Suhrawardi, the modes of light
and darkness contribute to the
ontological articulation of reality.
In Book I, Manifestation is
not treated as a linguistic
or cognitive representation
of an otherwise complete object.
Both approaches therefore
raise questions about the
ontological conditions under
which an existent appears
or becomes manifest.

% Structural ID: VS.B01.P01.C08.S14.SS16.P003
\phantomsection\label{para:VS.B01.P01.C08.S14.SS16.P003}

The first decisive divergence concerns
the role assigned to Light.
Suhrawardi makes light a central
ontological principle through
which the hierarchy and
differentiation of existents
are explained. Book I has not
adopted this account as the
definition of Light. Its
ontological status remains
subject to a separate
investigation. Consequently,
the two accounts cannot
be treated as equivalent
theories of Light.

% Structural ID: VS.B01.P01.C08.S14.SS16.P004
\phantomsection\label{para:VS.B01.P01.C08.S14.SS16.P004}

The second divergence concerns
the role of self-awareness.
Suhrawardi connects the
self-manifestation of pure
lights with their mode of
existence and with knowledge
by presence. Book I does
not make the existence of
Identity dependent upon
self-awareness or cognition.
This is not a denial that
self-aware entities exist.
It is a refusal to make
self-awareness a universal
condition of ontological
identity.

% Structural ID: VS.B01.P01.C08.S14.SS16.P005
\phantomsection\label{para:VS.B01.P01.C08.S14.SS16.P005}

The third divergence concerns
the status of the additional
domain of suspended images.
Suhrawardi's account and
its later commentarial
development assign a
distinct status to forms
that are neither ordinary
material bodies nor pure
intelligible lights. Book I
has not established that
such forms belong to one
of its proposed realms or
hubs. The correct result
is an open comparative
problem, not an immediate
assimilation.

% -----------------------------------------------------------------------------
% SS17 — Strongest Suhrawardian Challenge
% -----------------------------------------------------------------------------

\subsection{The Strongest Suhrawardian Challenge to Book I}
\label{sec:ch08-s14-ss17}

% Structural ID: VS.B01.P01.C08.S14.SS17.P001
\phantomsection\label{para:VS.B01.P01.C08.S14.SS17.P001}

The strongest Suhrawardian challenge concerns
the possibility that the conceptual
distinctions of Book I remain
insufficiently connected with
the ontological reality they
purport to articulate. If
light is self-evident and
if the reality of pure lights
is inseparable from their
self-manifestation, then an
ontology that begins with
analytical distinctions
must explain how those
distinctions correspond
to reality rather than
merely organize thought
about reality.

% Structural ID: VS.B01.P01.C08.S14.SS17.P002
\phantomsection\label{para:VS.B01.P01.C08.S14.SS17.P002}

A second form of the challenge
concerns the distinction
between Being and Identity.
Suhrawardi's critique of
the Peripatetic distinction
between essence and existence
may be used to question
whether the distinction
adopted in Book I is
ontologically meaningful.
If a distinction between
two concepts does not
correspond to two
independently existing
constituents, why should
the distinction between
Being and Identity be
regarded as more than
a conceptual division?

% Structural ID: VS.B01.P01.C08.S14.SS17.P003
\phantomsection\label{para:VS.B01.P01.C08.S14.SS17.P003}

A third form of the challenge
concerns the explanatory
adequacy of Identity.
Suhrawardi's ontology
offers a detailed account
of how different modes
of light and darkness
occupy different positions
within a hierarchy and
how relations among them
contribute to the
differentiation of reality.
If Book I asserts that
Identity is irreducible
to determination, structure,
manifestation, or relation,
it must explain what
positive ontological work
Identity performs beyond
the exclusions of these
concepts.

% Structural ID: VS.B01.P01.C08.S14.SS17.P004
\phantomsection\label{para:VS.B01.P01.C08.S14.SS17.P004}

The strongest challenge is therefore
not that Suhrawardi has
already disproved the
Being--Identity distinction.
It is that Book I must
show why its distinction
is ontologically necessary,
how it avoids becoming
a merely formal separation,
and how it explains the
identity of an existent
without reducing that
identity to the qualities,
relations, or degrees
through which the existent
is manifested.

% -----------------------------------------------------------------------------
% SS18 — Book I Challenge and Possible Suhrawardian Reply
% -----------------------------------------------------------------------------

\subsection{The Strongest Book I Challenge and Suhrawardi's Possible Reply}
\label{sec:ch08-s14-ss18}

% Structural ID: VS.B01.P01.C08.S14.SS18.P001
\phantomsection\label{para:VS.B01.P01.C08.S14.SS18.P001}

The strongest challenge from Book I concerns
the explanatory transition from
the hierarchy of light to the
identity of individual existents.
A hierarchy may explain dependence,
rank, intensity, and causal order.
It does not follow that it
explains what makes each
existent the identity that
it is. The same caution
applies to the distinction
between pure and accidental
lights: a classification
of ontological modes is
not yet a complete theory
of identity.

% Structural ID: VS.B01.P01.C08.S14.SS18.P002
\phantomsection\label{para:VS.B01.P01.C08.S14.SS18.P002}

Book I also asks whether the account
of self-manifestation can
explain entities whose
existence does not involve
self-awareness. If the
self-awareness of pure
lights is a defining
feature of their mode
of existence, the account
must still explain the
ontological status of
entities that do not
possess this feature.
The identity of a corporeal
barrier, for example,
cannot be dismissed merely
because it does not
manifest itself to itself
in the manner of a pure
light.

% Structural ID: VS.B01.P01.C08.S14.SS18.P003
\phantomsection\label{para:VS.B01.P01.C08.S14.SS18.P003}

A possible Suhrawardian reply would be that
the objection presupposes
an overly narrow account
of manifestation. The
hierarchy of lights does
not merely classify
entities according to
their degree of awareness.
It explains the different
ways in which entities
depend upon, receive,
or exercise manifestation.
The corporeal world is
not excluded from reality;
it occupies a lower and
different mode within
the same ontological order.
The reply would therefore
deny that the absence
of self-awareness entails
the absence of ontological
determination.

% Structural ID: VS.B01.P01.C08.S14.SS18.P004
\phantomsection\label{para:VS.B01.P01.C08.S14.SS18.P004}

This possible reply has force,
but it does not settle
the problem of Identity.
It explains how different
entities may occupy
different modes of
manifestation and
dependence. It does
not yet establish whether
Identity is reducible
to the mode of manifestation
or whether an additional
ontological distinction
is required. The dispute
therefore remains open
between an ontology
that gives explanatory
priority to the hierarchy
of light and an ontology
that distinguishes Being
and Identity without
assigning ontological
priority to either.

% -----------------------------------------------------------------------------
% SS19 — Bilateral Costs and Substitution Test
% -----------------------------------------------------------------------------

\subsection{Bilateral Costs and the Substitution Test}
\label{sec:ch08-s14-ss19}

% Structural ID: VS.B01.P01.C08.S14.SS19.P001
\phantomsection\label{para:VS.B01.P01.C08.S14.SS19.P001}

If Book I were to adopt Suhrawardi's ontology
of light without qualification,
it would have to accept a
substantial explanatory commitment:
the differentiation of existent
things would be grounded in
the hierarchy and modes of
light, and the account of
manifestation would be
organized through the
distinctions among pure
lights, accidental lights,
dark barriers, and dark
accidents. Such an adoption
would require Book I to
reconsider the status of
Light, the relation between
Identity and self-manifestation,
and the status of the
additional domain of
suspended images.

% Structural ID: VS.B01.P01.C08.S14.SS19.P002
\phantomsection\label{para:VS.B01.P01.C08.S14.SS19.P002}

If Suhrawardi's system were to adopt
the foundational commitments
of Book I, it would have
to explain why the identity
of an existent cannot be
exhausted by its mode of
luminosity, causal position,
or self-manifestation.
It would also have to
distinguish the analytical
difference between Being
and Identity from a
distinction between
independently existing
constituents. Such a
reconstruction would
alter the explanatory
role assigned to the
hierarchy of light.

% Structural ID: VS.B01.P01.C08.S14.SS19.P003
\phantomsection\label{para:VS.B01.P01.C08.S14.SS19.P003}

The substitution test confirms that
the central terms of the
two systems cannot be
replaced without changing
their meaning. ``Light
of Lights'' cannot be
substituted for ``First
Determination''; ``pure
light'' cannot be
substituted for ``Identity'';
``self-manifestation''
cannot be substituted
for ``Manifestation'' in
the technical sense of
Book I; and ``suspended
image'' cannot be
substituted for a realm
or hub of the present
framework. The relation
between the systems is
therefore neither
equivalence nor identity
of doctrine.

% Structural ID: VS.B01.P01.C08.S14.SS19.P004
\phantomsection\label{para:VS.B01.P01.C08.S14.SS19.P004}

The most defensible relation is a qualified
structural comparison.
Both accounts investigate
the ontological differentiation
of existent things and
the relation between
manifestation and reality.
They differ, however,
in the principles through
which they explain that
differentiation. Suhrawardi
organizes his account
around light, darkness,
and the hierarchy of
immaterial and material
entities. Book I preserves
the distinction between
Being and Identity and
does not adopt the
hierarchy of light as
its foundational ontology.

% -----------------------------------------------------------------------------
% SS20 — Open Problems, Claim Status, and Transition
% -----------------------------------------------------------------------------

\subsection{Open Problems, Claim Status, and Transition}
\label{sec:ch08-s14-ss20}

% Structural ID: VS.B01.P01.C08.S14.SS20.P001
\phantomsection\label{para:VS.B01.P01.C08.S14.SS20.P001}

The comparison leaves several problems open.
First, the precise ontological relation
between light and the identity of
individual existents requires further
analysis. Second, the distinction
between the original doctrine of
suspended images and its later
commentarial elaboration must be
preserved in any detailed account
of the imaginal domain. Third,
the relation between self-manifestation
and self-awareness must be examined
without making cognition a universal
condition of ontological identity.
Fourth, the ontological status of
Light within Book I remains
unresolved by the present
comparative inquiry.

% Structural ID: VS.B01.P01.C08.S14.SS20.P002
\phantomsection\label{para:VS.B01.P01.C08.S14.SS20.P002}

The functional and epistemic classifications
of the claims in this section
must remain distinct. Historical
reconstructions of Suhrawardi's
doctrine are not foundational
commitments of Book I.
The distinctions adopted
for comparative purposes
are methodological constructions.
The challenges formulated
against Book I are arguments
to be examined, not
established refutations.
The unresolved relation
between Identity, light,
and self-manifestation
remains an open problem.

% Structural ID: VS.B01.P01.C08.S14.SS20.P003
\phantomsection\label{para:VS.B01.P01.C08.S14.SS20.P003}

The following table records the principal
claims and their comparative
limits. It distinguishes
the historical reconstruction
from the functional role
of the claim and from
its epistemic or justificatory
status within the present
investigation.


% =============================================================================
% Project : THE PHILOSOPHY OF VIABILITY
% Book    : Book I — Ontology of Viability
% Book ID : VS.B01
% Chapter : Chapter 08 — Comparative Philosophical Examination
% Section : Section 14 — Suhrawardi
% Asset   : Table 01 — Claims and Comparative Limits
% Asset ID: VS.B01.P01.C08.S14.SS20.P003.TL01
% File    : 4_ASSETS/tables/VS.B01.P01.C08.S14.SS20.P003.TL01.suhrawardi_claims_and_comparative_limits.tex
%
% Purpose:
%   Record the principal historical reconstructions, their functional
%   forms, their epistemic status, and their limits for Book I.
% =============================================================================

\begin{table}[htbp]
\centering
\small
\renewcommand{\arraystretch}{1.25}

\begin{tabular}{
    p{0.15\linewidth}
    p{0.29\linewidth}
    p{0.22\linewidth}
    p{0.26\linewidth}
}
\toprule
\textbf{Issue}
&
\textbf{Suhrawardian Reconstruction}
&
\textbf{Claim Form and Status}
&
\textbf{Limit for Book I}
\\
\midrule

Light
&
Light is an ontologically central and self-evident principle
whose modes articulate the order of reality.
&
Form: Thesis.
Status: Historical reconstruction.
&
Does not establish that Book I Light is identical with
Suhrawardi's light.
\\

Pure lights
&
Pure immaterial lights are self-subsistent, incorporeal,
and manifest to themselves.
&
Form: Ontological proposition.
Status: Historical reconstruction.
&
Self-manifestation does not establish that Identity
is self-awareness.
\\

Dark barriers and accidents
&
Corporeal barriers and dependent dark accidents occupy
distinct positions within the ontology.
&
Form: Classification.
Status: Historical reconstruction.
&
Darkness must not be equated with non-being or with
the absence of Identity.
\\

Light of Lights
&
The first principle grounds the causal and ontological
order of dependent lights.
&
Form: Foundational commitment within Suhrawardi's system.
Status: Historical reconstruction of his argument.
&
Must not be equated with Being, Identity, or
First Determination.
\\

Vertical and horizontal orders
&
The hierarchy includes vertical dependence and
horizontal relations among lights of comparable rank.
&
Form: Ontological thesis.
Status: Historical reconstruction.
&
Does not map onto the four realms or three hubs
of Book I.
\\

Essence and existence
&
Existence and related notions are treated as
mental considerations in the relevant critique.
&
Form: Philosophical thesis.
Status: Interpretively disputed.
&
Not equivalent to the Being--Identity distinction.
\\

Suspended images
&
Suspended forms possess a mode of existence not
reducible to ordinary material inherence.
&
Form: Ontological thesis.
Status: Historical reconstruction with open
commentarial questions.
&
Does not establish identity with Relational
Manifestation or the Hub of Signification.
\\

Knowledge by presence
&
Self-awareness and unmediated presence are central
to Suhrawardi's account of knowledge.
&
Form: Epistemological thesis.
Status: Historical reconstruction.
&
Not adopted as an ontological condition of
Identity in Book I.
\\

Identity and individuation
&
Intensity, rank, relations, and accidents contribute
to the differentiation of lights.
&
Form: Comparative inference.
Status: Open problem.
&
No demonstrated reduction of Identity to
luminosity, intensity, or individuation.
\\

\bottomrule
\end{tabular}

\caption{Principal Suhrawardian claims and the limits of their
comparison with the ontological commitments of Book I.
Functional form and epistemic status are recorded separately
in the third column.}
\label{tab:ch08-s14-suhrawardi-claims}

\end{table}

% =============================================================================
% END OF TABLE 01 — SUHRAWARDI
% =============================================================================


% Structural ID: VS.B01.P01.C08.S14.SS20.P004
\phantomsection\label{para:VS.B01.P01.C08.S14.SS20.P004}

The examination of Suhrawardi has therefore
not established that Light
is Identity, that the
Light of Lights is
First Determination, or
that the suspended images
belong to one of the
proposed realms or hubs.
It has instead clarified
the ontological problems
that any account of
Identity and Manifestation
must confront: the
relation between self-subsistence
and dependence, the
differentiation of modes
of manifestation, the
status of entities that
are not self-aware, and
the possibility of
ontological domains
not captured by the
initial classification.

% Structural ID: VS.B01.P01.C08.S14.SS20.P005
\phantomsection\label{para:VS.B01.P01.C08.S14.SS20.P005}

The next section turns to Averroes.
The comparison will examine
his treatment of substance,
form, matter, causation,
intellect, and the relation
between philosophical
demonstration and the
structure of the existent.
The transition must preserve
the distinction between
Suhrawardi's ontology
of light and Averroes'
critical engagement with
the Aristotelian tradition.
Neither position will be
translated into the
vocabulary of Book I
without an independent
argument.

% =============================================================================
% END OF SECTION 14 OF CHAPTER 08
% =============================================================================
